% Project Strawberry 
\message{\optexversion}
\fontfam[LMfonts]
\typosize[11/13]
\margins/1 a4 (2.5,2.5,2.5,2.5)cm
\input definitions

\tit Strawberry - Phase 1

\chap Introduction

This project is about building an app that plays local audio (music) files. My goals will be listed in depth later but basically I am aiming for a minimal GUI with typical music-playing functionalities (e.g. shuffle, looping, volume etc.) and easily customisable organisation (e.g. albums, artists). I will be building this app specifically for my operating system of arch-linux. I have basically no experience with app development so this will be a very informal project and I will try to explain what I do and why. This project is titled phase 1 because if I finish this I will try to make a physical device to run the app and play music outside (an MP3 player). 

\sec Obtaining Music

The main ways of obtaining music files is either by digitising a CD, buying the music online (e.g. amazon music) or by downloading copyright-free music online. To avoid any copyright issues I have downloaded a selection of copyright-free classical music recordings using \code{yt-dlp} and \code{ffmpeg}. These packages make it easier to download and convert music. I decided to use opus file format because it takes less space and I'm building from scratch so there shouldn't be any compatibility issues. I downloaded the music using the terminal command `yt-dlp -x --audio-format opus --audio-quality 128K --ppa "ExtractAudio+ffmpeg_o:-acodec libopus -af loudnorm" <URL>`, where \code{-x} only downloads audio and \code{--ppa "..."} is an ffmpeg command to normalise audio (long because it forces recompilation to opus as yt-dlp preferes to copy a file if it exists which messes up the audio normalisation). I also messed around with downloading to specific files and auto-inserting metadata but decided it would be better to make a simple editor as part of the app.  

\bye
